% \textbf{Step 1: Abstract Transition Graph}
\textbf{Step 1: }
The Section~\ref{sec:psRB-progabs} first 
computes the \emph{Abstract Transition Graph} as in Figure~\ref{fig:relatedNestedWhileOdd-overview}(b).
Each edge $l \xrightarrow{dc} l'$ is an abstract transition $\absevent = (l, dc, l')$ annotated with a constraint $dc$ corresponding to the command of label $l$.

% \textbf{Step 2: Program Refinement}
\textbf{Step 2: }
The second step in Section~\ref{sec:psRB-refine}
computes the \emph{Refined Program}, $\rprog$ for a program $c$ based on 
its abstraction transition graph and transforms the multiple-paths loops
into multiple loops where
the interleaving of paths is explicit as in bottom part of Figure~\ref{fig:relatedNestedWhileOdd-overview}(c).
It has two interleaving patterns in the loop $L_1$.
%  denoted as $\rprog_1^1$ and $\rprog_1^2$.

% \textbf{Step 3: Ranking Function Estimation}
\textbf{Step 3: }
In the meanwhile, Section~\ref{sec:psRB-rank} computes the \emph{Ranking Function}, $\locbound(\absevent, c)$ 
for every edge $\absevent$ 
and estimates an upper bound invariant for each.

% \textbf{Step 4: Path-sensitive Reachability-bound Computation.}
\textbf{Step 4: }
% The \emph{path-sensitive reachability-bound} 
The algorithm computes the \emph{Reachability-bound}, $\psRB(l, c)$ for every program point $l$ using the $\rprog$ and the upper bound invariant of the $\locbound(\absevent, c)$ where $\absevent = (l,dc,l')$.
It requires to compute our two novel quantities, the \emph{Path Reachability-bound}, $\inoutB(\rprog, \tpath, c)$ and the \emph{Loop Reachability-bound}, $\lpchB(l: \rprog, \tpath, c)$.
Section~\ref{sec:psRB-psrb} introduces this algorithm and the following sections describe the computations. 
The soundness of each algorithm is in the Appendices and the input $c$ will be omitted if the context is clear.

In the first interleaving pattern, $\rprog_1^1 = 1: \rprepeat(\tpath_1; 4:\rprepeat(\tpath_3); \tpath_2; \tpath_4)$ in Figure~\ref{fig:relatedNestedWhileOdd-overview},
we first compute $\outinB(4:\rprepeat(\tpath_3), \tpath_3) = n - m$
 for $\tpath_3$ in its innermost loop $L_4$ as a local \emph{path reachability-bound} in Section~\ref{sec:psRB-pathlocalrb}.
Then we compute $\lpchB(\rprog_1^1, \tpath_3) = 1$ w.r.t. its outer loop $L_1$ in Section~\ref{sec:psRB-looprb}. In the second interleaving pattern, we compute $\outinB(4:\rprepeat(\tpath_3), \tpath) = n - m - 3$ and the same for $\lpchB$.
So Section~\ref{sec:psRB-pathrb} computes $\inoutB(\rprog, \tpath_3) = \max\{ 1 \times (n - m), 1 \times (n - m - 3) \} = n - m$ globally.

Then for every program point $l$, we sum up all the $\inoutB(\rprog, \tpath)$ over $\tpath$ that contains $l$ and get $\psRB(l, c)$.
Since point $5$ only shows up on $\tpath_3$, we compute \highlight{$\psRB(5, c) = n - m$}.
The points $0$ and $\lex$ are not in any loop, so $\psRB(0) = \psRB(\lex) = 1$,
The points $3, 6, 7$ and $8$ which only show up once on $\tpath_2$ and $\tpath_4$ are all equal to $\lfloor\frac{m}{4}\rfloor$, which are same as their $\inoutB$.
For the loop headers $1$ and $4$, we only sum up the $\inoutB(\rprog, \tpath)$ where they show up as start-point of the $\tpath$.
So $\psRB(4) =  \lfloor\frac{m}{4}\rfloor + n - m + 1$ and $\psRB(1) = 2 \times \lfloor\frac{m}{4}\rfloor + 1$ all as expected.
